% capitolo 3: analisi dei requisiti
\chapter{Analisi dei requisiti}
\label{cap:analisi-requisiti}
% TENERE SEMPRE A MENTE CHE OGNI REQUISITO DEVE RISPETTARE QUESTI 2 PRINCIPI:
% 1. Necessità: ogni requisito inserito nella lista deve soddisfare un bisogno aziendale reale ed esplicito.
% 2. Sufficienza: I requisiti elencati devono coprire tutti i bisogni dell'azienda.

 \section{Introduzione}
 In questa sezione vengono analizzati i casi d'uso e i requisiti aziendali necessari per l'introduzione di un \textit{Data Catalog} nell'infrastruttura di Wintech. 
 Dato che le moderne soluzioni software offrono un ampio spettro di funzionalità, non sempre necessarie, l'analisi ha lo scopo di filtrare gli strumenti attualmente 
 presenti sul mercato, individuando quelli che meglio rispondono ai requisiti dell'azienda e alle sue reali esigenze operative.

\subsection{Attori del sistema}
In ambito aziendale possono essere molteplici le figure che hanno la necessità di accedere a un \textit{Data Catalog}, ognuna con esigenze e competenze diverse. 
Ai fini dell'analisi dei requisiti, si individuano i seguenti attori principali:
\begin{itemize}
    \item \textbf{Analista o Sviluppatore di BI}: Utente finale del catalogo, il cui obiettivo è ricercare e accedere ai dati provenienti da varie fonti, comprenderne il significato e 
    utilizzarli per generare report e analisi o per esportarli in strumenti esterni di Business Intelligence.
    \item \textbf{Responsabile della Governance e della sicurezza dei dati}: Figura garante della qualità e della sicurezza dei dati, che controlla che i dati siano corretti e che rispettino le normative vigenti.
    \item \textbf{Amministratore di sistema}: Figura tecnica che gestisce l'infrastruttura della piattaforma e connette le varie fonti di dati aziendali e mantiene il catalogo aggiornato, garantendone il corretto funzionamento.
\end{itemize}

\newpage

\section{Casi d'uso}
In questa sezione vengono analizzati i casi d'uso e gli scenari per cui l'azienda ha bisogno di adottare un \textit{Data Catalog}. 
Per maggiore chiarezza organizzativa, i casi d'uso vengono suddivisi in tre macroaree raggruppate in base all'attore primario coinvolto.\\
Ogni caso d'uso è identificato dalla seguente notazione:
\begin{center}
    \textbf{UC-[Macroarea].[Numero]}
\end{center}

% MACROAREA 1
\subsection{Amministrazione del sistema}

\begin{usecase}{1.0}{Collegamento a una fonte dati strutturata}
    \usecaseactors{Amministratore di sistema}
    \usecasedesc{L'amministratore collega una fonte dati strutturata al catalogo dati aziendale e ne visualizza i metadati tramite interfaccia}
    \usecasepre{L'attore ha effettuato l'accesso alla piattaforma con privilegi di amministratore e possiede le credenziali della fonte dati da importare}
    \usecasepost{I metadati della fonte dati selezionata sono stati importati con successo e sono visibili all'interno del catalogo dati}
    \usecasemain{
        \begin{enumerate}
            \item L'attore seleziona l'opzione che permette di connettere una nuova fonte di dati al catalogo.
            \item L'attore seleziona la fonte desiderata e inserisce credenziali o parametri per la connessione.
            \item Il sistema verifica la connessione e notifica il successo della connessione.
            \item L'attore avvia l'estrazione dei metadati.
        \end{enumerate}
    }
    \usecasealt{
        \begin{enumerate}
            \item Il sistema segnala il fallimento della connessione.
            \item L'attore non può importare i dati nel catalogo e deve ricontrollare i parametri per la connessione.
        \end{enumerate}
    }
    \label{uc:1_0_fonte_strutturata}
\end{usecase}

\begin{usecase}{1.1}{Collegamento a una fonte dati semi strutturata}
    \usecaseactors{Amministratore di sistema}
    \usecasedesc{L'amministratore collega una fonte dati semi strutturata al catalogo dati aziendale e ne visualizza i metadati tramite interfaccia}
    \usecasepre{L'attore ha effettuato l'accesso alla piattaforma con privilegi di amministratore e possiede le credenziali o il percorso della fonte dati da importare}
    \usecasepost{I metadati della fonte dati selezionata sono stati importati con successo e sono visibili all'interno del catalogo dati}
    \usecasemain{
        \begin{enumerate}
            \item L'attore seleziona l'opzione che permette di aggiungere una nuova fonte di dati al catalogo.
            \item L'attore seleziona la fonte desiderata e inserisce i parametri per la connessione o il percorso del file.
            \item Il sistema verifica la connessione e notifica il successo della connessione.
            \item L'attore avvia l'estrazione dei metadati.
        \end{enumerate}
    }
    \label{uc:1_1_fonte_semi_strutturata}
\end{usecase}

\begin{usecase}{1.2}{Collegamento a una fonte dati non strutturata}
    \usecaseactors{Amministratore di sistema}
    \usecasedesc{L'amministratore collega una fonte dati non strutturata al catalogo dati aziendale e ne visualizza i metadati tramite l'interfaccia della piattaforma}
    \usecasepre{L'attore ha effettuato l'accesso alla piattaforma con privilegi di amministratore e possiede le credenziali o il percorso della fonte dati da importare}
    \usecasepost{I metadati della fonte dati selezionata sono stati importati con successo e sono visibili all'interno del catalogo dati}
    \usecasemain{
        \begin{enumerate}
            \item L'attore seleziona l'opzione che permette di aggiungere una nuova fonte di dati al catalogo.
            \item L'attore seleziona la fonte desiderata e inserisce i parametri per la connessione o il percorso del file.
            \item Il sistema verifica la connessione e notifica il successo della connessione.
            \item L'attore avvia l'estrazione dei metadati.
        \end{enumerate}
    }
    \label{uc:1_2_fonte_non_strutturata}
\end{usecase}

\begin{usecase}{1.3}{Visualizzazione delle fonti dati collegate}
    \usecaseactors{Amministratore di sistema}
    \usecasedesc{L'amministratore visualizza l'inventario delle fonti dati, cioè l'insieme delle fonti di dati che popolano il \textit{Data Catalog}}
    \usecasepre{L'attore ha effettuato l'accesso alla piattaforma con privilegi di amministratore}
    \usecasepost{L'attore accede all'elenco completo delle fonti di dati importate nel \textit{Data Catalog}}
    \usecasemain{
        \begin{enumerate}
            \item L'attore seleziona l'opzione per visualizzare l'elenco delle fonti dati del catalogo dati.
            \item Il sistema mostra l'elenco completo delle fonti dati.
        \end{enumerate}
    }
    \label{uc:1_3_visualizzazione_fonti_collegate}
\end{usecase}

\begin{usecase}{1.4}{Gestione dei permessi e degli accessi}
    \usecaseactors{Amministratore di sistema}
    \usecasedesc{L'amministratore definisce i ruoli degli utenti che accedono alla piattaforma e assegna i permessi di accesso ai diversi set di dati. 
    Il sistema garantisce che ogni utente visualizzi esclusivamente i dati per cui dispone delle autorizzazioni necessarie}    
    \usecasepre{L'attore ha effettuato l'accesso alla piattaforma con privilegi di amministratore}
    \usecasepost{I ruoli degli utenti sono stati definiti e i permessi di accesso ai dati sono stati configurati correttamente}
    \usecasemain{
        \begin{enumerate}
            \item L'attore seleziona l'opzione per la gestione dei ruoli e dei permessi.
            \item L'attore crea o modifica i ruoli disponibili nella piattaforma.
            \item L'attore assegna i ruoli agli utenti che accedono al catalogo.
            \item L'attore definisce le regole di accesso ai dati associate a ciascun ruolo.
            \item Il sistema applica i permessi configurati e limita la visibilità dei dati in base al ruolo di ciascun utente.
        \end{enumerate}
    }
    \label{uc:1_4_gestione_permessi}
\end{usecase}

\begin{usecase}{1.5}{Pianificazione della sincronizzazione dei metadati}
    \usecaseactors{Amministratore di sistema}
    \usecasedesc{L'amministratore pianifica la sincronizzazione automatica dei metadati, definendo la frequenza e le modalità con cui il catalogo si aggiorna allineandosi 
    alle fonti dati originali, garantendo che le informazioni presenti nel catalogo siano sempre coerenti e aggiornate}
    \usecasepre{L'attore ha effettuato l'accesso alla piattaforma con privilegi di amministratore e almeno una fonte dati è stata collegata al catalogo}
    \usecasepost{La sincronizzazione automatica è stata pianificata correttamente e i metadati vengono aggiornati secondo la frequenza stabilita}
    \usecasemain{
        \begin{enumerate}
            \item L'attore seleziona le impostazioni per la pianificazione della sincronizzazione dei metadati.
            \item L'attore configura la frequenza e le modalità di aggiornamento desiderate.
            \item Il sistema conferma la configurazione e pianifica le sincronizzazioni secondo i parametri definiti.
            \item Il sistema esegue automaticamente la sincronizzazione dei metadati agli intervalli stabiliti.
        \end{enumerate}
    }
    \label{uc:1_5_pianificazione_sincronizzazione}
\end{usecase}

\begin{usecase}{1.6}{Configurazione della classificazione automatica dei dati sensibili}
    \usecaseactors{Amministratore di sistema}
    \usecasedesc{L'amministratore abilita e configura il motore di classificazione automatica dei dati sensibili, che analizza i metadati delle fonti dati collegate 
    al catalogo e individua in autonomia i campi contenenti dati personali o sensibili (PII). Il motore si avvale per esempio di tecniche di Intelligenza Artificiale e di espressioni 
    regolari per riconoscere pattern tipici di dati sensibili, quali nomi, indirizzi email, numeri di telefono, codici fiscali e altri dati identificativi}
    \usecasepre{L'attore ha effettuato l'accesso alla piattaforma con privilegi di amministratore e almeno una fonte dati è stata collegata al catalogo}
    \usecasepost{Il motore di classificazione automatica è attivo e configurato. Durante le sincronizzazioni dei metadati, il sistema individua e contrassegna automaticamente i campi contenenti dati sensibili}
    \usecasemain{
        \begin{enumerate}
            \item L'attore accede alle impostazioni della classificazione automatica dei dati sensibili.
            \item L'attore abilita il motore di classificazione e seleziona le fonti dati su cui applicarlo.
            \item L'attore configura le categorie di dati sensibili da rilevare o adotta la configurazione predefinita della piattaforma.
            \item Il sistema conferma la configurazione e attiva il motore.
            \item Ad ogni sincronizzazione, il sistema analizza i metadati delle fonti selezionate e contrassegna automaticamente i campi identificati come sensibili.
        \end{enumerate}
    }
    \label{uc:1_6_configurazione_classificazione_dati}
\end{usecase}

\begin{usecase}{1.7}{Monitoraggio dello stato delle sincronizzazioni e delle connessioni}
    \usecaseactors{Amministratore di sistema}
    \usecasedesc{L'amministratore visualizza una dashboard tecnica, un report o un avviso che indica lo stato delle connessioni alle fonti dati e l'esito 
    delle sincronizzazioni dei metadati, consentendo di verificare che le operazioni pianificate siano avvenute correttamente e di individuare eventuali anomalie o errori}
    \usecasepre{L'attore ha effettuato l'accesso alla piattaforma con privilegi di amministratore e almeno una sincronizzazione è stata pianificata ed eseguita}
    \usecasepost{L'attore ha visualizzato lo stato aggiornato delle connessioni e delle sincronizzazioni ed è a conoscenza di eventuali errori o anomalie da risolvere}
    \usecasemain{
        \begin{enumerate}
            \item L'amministratore accede alla dashboard di monitoraggio delle connessioni e delle sincronizzazioni.
            \item Il sistema mostra lo stato di ciascuna operazione, indicando se è avvenuta con successo o se si sono verificati errori.
            \item In caso di errore, il sistema fornisce informazioni dettagliate sulla causa del problema, ad esempio una credenziale scaduta o una fonte dati non raggiungibile.
            \item L'attore prende visione delle anomalie e interviene per risolvere il problema alla fonte.
        \end{enumerate}
    }
    \label{uc:1_7_monitoraggio_sincronizzazioni}
\end{usecase}

\begin{usecase}{1.8}{Gestione di avvisi e alert tecnici}
    \usecaseactors{Amministratore di sistema}
    \usecasedesc{Il sistema invia automaticamente alert di natura tecnica e avvisi all'amministratore quando vengono rilevate anomalie tecniche, come il fallimento di una sincronizzazione, 
    un errore di connessione o la scadenza delle credenziali. L'amministratore visualizza gli alert ricevuti e interviene per risolvere i problemi segnalati}
    \usecasepre{L'attore ha effettuato l'accesso alla piattaforma con privilegi di amministratore e il sistema di notifiche automatiche è stato configurato}
    \usecasepost{L'attore ha preso visione degli alert ricevuti e ha avviato le azioni correttive necessarie per risolvere i problemi segnalati}
    \usecasemain{
        \begin{enumerate}
            \item Il sistema rileva un'anomalia e genera automaticamente un alert.
            \item L'attore riceve una notifica e accede alla sezione dedicata agli alert e agli avvisi.
            \item Il sistema mostra i dettagli dell'alert, come la causa, la fonte dati coinvolta e il momento in cui si è verificato il problema.
            \item L'attore esamina le informazioni disponibili e interviene per risolvere il problema.
            \item Il sistema non mostra più l'alert o lo marchia come "risolto".
        \end{enumerate}
    }
    \label{uc:1_8_gestione_alert_tecnici}
\end{usecase}


% MACROAREA 2
\subsection{Governance e sicurezza dei dati}
\begin{usecase}{2.0}{Individuazione automatica dei dati sensibili}
    \usecaseactors{Responsabile della Governance e della sicurezza dei dati}
    \usecasedesc{La piattaforma analizza automaticamente i metadati dei dataset caricati nel catalogo e individua i campi contenenti dati sensibili o personali (PII), classificandoli tramite tag appositi. 
    Il processo avviene in modo trasparente durante la sincronizzazione dei metadati, senza richiedere intervento manuale da parte dell'attore}
    \usecasepre{L'attore ha effettuato l'accesso alla piattaforma con privilegi di responsabile della Governance e il motore di classificazione automatica è già stato configurato dall'amministratore}
    \usecasepost{I campi contenenti dati sensibili o personali sono stati identificati e classificati automaticamente con i tag appropriati}
    \usecasemain{
        \begin{enumerate}
            \item La piattaforma rileva l'aggiunta o la sincronizzazione di un nuovo dataset nel catalogo.
            \item Il motore di classificazione automatica analizza i metadati del dataset e individua i campi che corrispondono a pattern tipici di dati sensibili o personali.
            \item Il sistema associa automaticamente i tag appropriati ai campi individuati.
            \item L'attore riceve una notifica dell'avvenuta classificazione automatica e può procedere alla revisione o all'approvazione.
        \end{enumerate}
    }
    \label{uc:2_0_individuazione_automatica_dati}
\end{usecase}




% MACROAREA 3
\subsection{Analisi dei dati e Business Intelligence}


\newpage

\section{Requisiti}
In questa sezione vengono elencati e descritti i requisiti derivanti dallo studio dei casi d'uso, con l'obiettivo di avere un quadro chiaro delle caratteristiche 
che un \textit{Data Catalog} deve avere per rispondere alle esigenze aziendali. Per ogni requisito vengono forniti una descrizione e il tracciamento verso i casi d'uso correlati.\\

Ogni requisito è identificato dalla seguente notazione:
\begin{center}
    \textbf{R[Tipo]-[Importanza]-[Numero]}
\end{center}
dove le voci assumono i seguenti significati:
\begin{description}
    \item[Tipo] Indica la natura architetturale del requisito:
    \begin{itemize}
        \item \textbf{F (Funzionale)}: Descrive i servizi e le funzionalità tangibili che il sistema deve fornire.
        \item \textbf{V (Vincolo)}: Descrive limitazioni tecnologiche o normative imposte esternamente.
        \item \textbf{Q (Qualità)}: Descrive le caratteristiche qualitative misurabili del sistema.
    \end{itemize}

    \item[Importanza] Indica il grado di priorità per l'azienda:
    \begin{itemize}
        \item \textbf{O (Obbligatorio)}: Requisito irrinunciabile; la sua assenza compromette gli obiettivi primari del progetto.
        \item \textbf{D (Desiderabile)}: Requisito non strettamente vincolante nel breve termine, ma a valore aggiunto riconoscibile.
        \item \textbf{OP (Opzionale)}: Requisito di priorità minore, la cui implementazione può essere rimandata.
    \end{itemize}

    \item[Numero] Un intero progressivo univoco per ogni tipologia.
\end{description}

% TABELLE DEI REQUISITI
\subsection{Requisiti funzionali}
I requisiti funzionali descrivono i comportamenti e i servizi che il sistema deve fornire. Rispondono alla domanda: \textit{"Cosa deve fare il software?"}.


\subsection{Requisiti di vincolo}
I requisiti di vincolo rappresentano le limitazioni tecnologiche, normative, architetturali o implementative imposte dall'esterno entro le quali il progetto deve muoversi.

\subsection{Requisiti di qualità}
I requisiti di qualità descrivono le caratteristiche misurabili che il sistema deve garantire, come prestazioni, sicurezza, usabilità o manutenibilità.




\newpage

\section{Tracciamento dei requisiti}
% tabella di tracciamento fonte - requisito




% ROBA DEL TEMPLATE ==================================

% \begin{usecase}{0}{Scenario principale}
% \usecaseactors{Sviluppatore applicativi}
% \usecasepre{Lo sviluppatore è entrato nel plug-in di simulazione all'interno dell'IDE}
% \usecasedesc{La finestra di simulazione mette a disposizione i comandi per configurare, registrare o eseguire un test}
% \usecasepost{Il sistema è pronto per permettere una nuova interazione}
% \label{uc:scenario-principale}
% \end{usecase}


% \begin{table}%
% \caption{Tabella del tracciamento dei requisti funzionali}
% \label{tab:requisiti-funzionali}
% \begin{tabularx}{\textwidth}{lXl}
% \hline\hline
% \textbf{Requisito} & \textbf{Descrizione} & \textbf{Use Case}\\
% \hline
% RFN-1     & L'interfaccia permette di configurare il tipo di sonde del test & UC1 \\
% \hline
% \end{tabularx}
% \end{table}
